The paper~\cite{on_the_link_between_binomial_theorem_and_discrete_convolution} reveals the connection
between binomial theorem and discrete convolution of powers, via~\eqref{theorem:odd-power-identity}.
The works~\cite{glaisher19121n,neuman1977evaluation,carlitz1977note} are useful because
of discussion of convoluted power sums,
for instance the sums
\begin{align*}
    C(i,j;N) = \sum_{k=0}^{N} k^i (N-k)^j \quad \mathrm{and} \quad C(i,j;N) = \sum_{k=0}^{N} (k+a)^i (N-k-a)^j,
\end{align*}
which are used odd power identities
\begin{align*}
    n^{2m+1} = \sum_{r=0}^{m} \coeffA{m}{r} \sum_{k=1}^{n} k^r (n-k)^r
\end{align*}
and binomial form
\begin{align*}
(n+2a)^{2m+1} = \sum_{r=0}^{m} \coeffA{m}{r} \sum_{k=-a+1}^{n+a} (k+a)^r (n+a-k)^r.
\end{align*}
